\documentclass[11pt]{article}
\usepackage{amsmath}
\title{When seeing a formation gap requires a sustained view}
\begin{document}
\maketitle

The optimal point sample only touches the detection threshold.
It does not provide an interval of usable observations. A
detector that needs a sustained view changes the optimum, even
with zero additional transmission delay.

\paragraph{Law.}
Keep the forward-only two-robot model with \(a=V=1,A=4,W=3\),
time 4, target common rest-to-rest advance \(D\), and physical
gap band \([1,8/5]\). The hidden initial gaps are \(6/5\) and
\(7/5\). The detection range is \(6/5-d\), \(0<d\le3/800\).
Now the controller may schedule one interval \([s,r]\) of
length \(h>0\). It obtains a valid detection at \(r\) only if
the view was available for the whole interval. It gets no
partial readout before \(r\). A valid negative is distinct
from erasure. Assume \(0<h\le1/10\).
This is an authored acquisition law, not a guarantee from
a camera, optical receiver, or radio source.

\paragraph{Exact delay and advance.}
Write \(y=x_2-x_1\) and \(t_0=\sqrt{2d}\). While the prefix
is common, the far gap cannot be detected anywhere in the
safe intersection band. A useful view requires \(y(t)\le-d\)
throughout \([s,r]\); in particular \(s\ge t_0\). The exact
advance is \(D=143/40-\lambda\), where
\[
  \lambda=
  \begin{cases}
    \sqrt{3d/2+9h^2/16}, &h\le2t_0/3,\\
    f(t_0+h),&h\ge2t_0/3,
  \end{cases}
  \qquad f(t)=t-\frac12\sqrt{t^2-2d}.
\]
The two expressions meet at \(\lambda=t_0\). In the short case
\(\lambda\le t_0\); in the long case
\(\lambda\le t_0+h/2\), using \(\sqrt{h^2+2t_0h}\ge h\).
Thus \(\lambda<0.137\) and the asserted advance exceeds \(17/5\).

For the upper bound, consider a distinguishing interval ending
by time \(1/2\). At each \(t\) in it the point-sample argument
bounds robot 2's advance at time 1 by \(7/8-f(t)\).
That argument still applies up to \(1/2\): the cut speed bound
\(2\sqrt{t^2-2d}\le1\) is valid and at least half a unit of
time remains to reach speed one. It applies even though the
controller does not learn the observation until the end.
For the short-dwell case, write \(b\) for the loss from \(7/8\).
If \(b\ge t_0\), it is already at least \(\lambda\).
Otherwise \(f(t)\le b\) for \(t\) in the view interval, and
\(t-b>0\), so
\[
 f(t)\le b \quad\Longleftrightarrow\quad
         \frac32t^2-4bt+2b^2+d\le0 .
\]
The permitted times have width at most
\(\sqrt{16b^2/9-8d/3}\), implying \(b\ge\lambda\).
For a longer dwell, \(r\ge t_0+h\), which lies on the
increasing branch of \(f\) (its minimum is at
\(2t_0/\sqrt3\)). Thus \(b\ge f(r)\ge f(t_0+h)\).
For \(r\in[1/2,1]\) the cruder speed bound gives
\(x_2(1)\le5/8-d<7/8-\lambda\).
If there is no distinguishing read before time 1, the shared
prefix gives \(x_2(1)\le7/10\). In all cases the near-case
middle and final advances are at most \(2\) and \(7/10\).

\paragraph{Attaining the bound.}
For the short dwell choose \(s=4\lambda/3-h/2\),
\(r=4\lambda/3+h/2\). The bound on \(h\) gives
\(\lambda\le s\), and these two values are the roots of
the quadratic above with \(b=\lambda\); \(s\ge t_0\).
For the longer dwell start at \(s=t_0\) and end at
\(r=t_0+h\), with \(\lambda\ge s\).
In both cases the end is below \(1/4\) under the declared
parameter bounds. Robot 1 accelerates at one throughout the
first unit. Robot 2 waits until \(\lambda\) and accelerates
at four, capping its speed at one after the view interval.
Both endpoints have relative displacement \(-d\):
in the short case by the root equations; in the long case
by zero advance at \(s\) and the definition of \(f(r)\).
Inside the interval \(y\le-d\): before the ramp this follows
from \(t\ge t_0\); on its overlap with the ramp the quadratic
relative displacement is convex and at most its values at
the two ends. The most negative relative gap in the prefix
is \(-2\lambda^2/3\), safely inside the intersection band.

At delivery the speed difference \(w=3r-4\lambda\) is
nonnegative, since the chosen end lies on the rising side
of the relative-displacement trajectory. In the near
branch robot 2 continues its fast ramp to the cap. Its
relative lead at 1 is \(p=3/8-\lambda>1/5\).
For the far branch, a feasible slower option brakes robot 2
at four until its relative speed reaches zero, then follows
robot 1's initial velocity through time 1. Its final relative
gap is \(-d+w^2/10<1/5\), using \(w\le r<1/4\).
Mix this option with the fast ramp so the relative lead at 1
is \(1/5\). Both have nondecreasing relative gap after receipt.
In the middle and final intervals use the scaled triangular
speed adjustment of area \(p+1/5\) and reversed ramps with
end lags \(1/5\) and \(p\). They attain \(D\) at rest.
If the report is erased, use the far branch's initial option,
middle area \(2/5\), and end lag \(1/5\). This retains the
intersection band and achieves \(17/5\), not the larger target.

\paragraph{Check and scope.}
For \(d=63/20000\), \(h=1/25\), the exact interval is
\([2/25,3/25]\), with delay \(3/40\) and advance \(7/2\).
The replay checks continuous velocity and gap bounds for
five rational cases, including a longer dwell that must begin at the earliest geometrically possible time. A sampled two-world trajectory LP
can attain the same figure at this interval; it cannot
start the interval at 0.075 and gives lower advance if
the interval is moved to start at 0.085. The proof supplies
the exact optimum over continuous time; the LP checks only the
sampled comparator. Dwells outside the declared range, other sensor laws, speed
changes, clock error and real evidence delivery are not
covered. This is partial R3a/R6, not full contract acceptance.
\end{document}
